% Options for packages loaded elsewhere
\PassOptionsToPackage{unicode}{hyperref}
\PassOptionsToPackage{hyphens}{url}
\PassOptionsToPackage{dvipsnames,svgnames,x11names}{xcolor}
\documentclass[
  a4paper,
]{article}
\usepackage{xcolor}
\usepackage[margin=25mm]{geometry}
\usepackage{amsmath,amssymb}
\setcounter{secnumdepth}{-\maxdimen} % remove section numbering
\usepackage{iftex}
\ifPDFTeX
  \usepackage[T1]{fontenc}
  \usepackage[utf8]{inputenc}
  \usepackage{textcomp} % provide euro and other symbols
\else % if luatex or xetex
  \usepackage{unicode-math} % this also loads fontspec
  \defaultfontfeatures{Scale=MatchLowercase}
  \defaultfontfeatures[\rmfamily]{Ligatures=TeX,Scale=1}
\fi
\usepackage{lmodern}
\ifPDFTeX\else
  % xetex/luatex font selection
  \setmainfont[]{DejaVu Serif}
  \setsansfont[]{DejaVu Sans}
  \setmonofont[]{DejaVu Sans Mono}
\fi
% Use upquote if available, for straight quotes in verbatim environments
\IfFileExists{upquote.sty}{\usepackage{upquote}}{}
\IfFileExists{microtype.sty}{% use microtype if available
  \usepackage[]{microtype}
  \UseMicrotypeSet[protrusion]{basicmath} % disable protrusion for tt fonts
}{}
\makeatletter
\@ifundefined{KOMAClassName}{% if non-KOMA class
  \IfFileExists{parskip.sty}{%
    \usepackage{parskip}
  }{% else
    \setlength{\parindent}{0pt}
    \setlength{\parskip}{6pt plus 2pt minus 1pt}}
}{% if KOMA class
  \KOMAoptions{parskip=half}}
\makeatother
\usepackage{longtable,booktabs,array}
\newcounter{none} % for unnumbered tables
\usepackage{calc} % for calculating minipage widths
% Correct order of tables after \paragraph or \subparagraph
\usepackage{etoolbox}
\makeatletter
\patchcmd\longtable{\par}{\if@noskipsec\mbox{}\fi\par}{}{}
\makeatother
% Allow footnotes in longtable head/foot
\IfFileExists{footnotehyper.sty}{\usepackage{footnotehyper}}{\usepackage{footnote}}
\makesavenoteenv{longtable}
\usepackage{graphicx}
\makeatletter
\newsavebox\pandoc@box
\newcommand*\pandocbounded[1]{% scales image to fit in text height/width
  \sbox\pandoc@box{#1}%
  \Gscale@div\@tempa{\textheight}{\dimexpr\ht\pandoc@box+\dp\pandoc@box\relax}%
  \Gscale@div\@tempb{\linewidth}{\wd\pandoc@box}%
  \ifdim\@tempb\p@<\@tempa\p@\let\@tempa\@tempb\fi% select the smaller of both
  \ifdim\@tempa\p@<\p@\scalebox{\@tempa}{\usebox\pandoc@box}%
  \else\usebox{\pandoc@box}%
  \fi%
}
% Set default figure placement to htbp
\def\fps@figure{htbp}
\makeatother
\setlength{\emergencystretch}{3em} % prevent overfull lines
\providecommand{\tightlist}{%
  \setlength{\itemsep}{0pt}\setlength{\parskip}{0pt}}
\usepackage{luatexja}
\usepackage{xurl}
\usepackage{bookmark}
\IfFileExists{xurl.sty}{\usepackage{xurl}}{} % add URL line breaks if available
\urlstyle{same}
\hypersetup{
  pdftitle={The Thirteenth Thought Experiment: Why Does a Particle That Has Acquired Mass Not Dissipate and Vanish?},
  pdfauthor={Noriaki Kihara (WF System Co., Ltd.) ORCID 0009-0004-6753-4020},
  colorlinks=true,
  linkcolor={Maroon},
  filecolor={Maroon},
  citecolor={Blue},
  urlcolor={Blue},
  pdfcreator={LaTeX via pandoc}}

\title{The Thirteenth Thought Experiment: Why Does a Particle That Has
Acquired Mass Not Dissipate and Vanish?}
\usepackage{etoolbox}
\makeatletter
\providecommand{\subtitle}[1]{% add subtitle to \maketitle
  \apptocmd{\@title}{\par {\large #1 \par}}{}{}
}
\makeatother
\subtitle{― Coherence as a Relation, Harmonic Exchange of Localized
Waves, Internal Clocks and Particle Lifetimes Considered from a Single
Interaction ―}
\author{Noriaki Kihara (WF System Co., Ltd.) ORCID 0009-0004-6753-4020}
\date{October 10, 2026 Version DOI 10.5281/zenodo.23266251}

\ltjsetparameter{jacharrange={-2,-3}}
\begin{document}
\maketitle

\textbf{Author:} Noriaki Kihara\\
\textbf{ORCID:} 0009-0004-6753-4020\\
\textbf{Version DOI:} 10.5281/zenodo.23266251\\
\textbf{Concept DOI:} 10.5281/zenodo.23266250\\
\textbf{Version:} v1.0\\
\textbf{Date:} 2026-10-10\\
\textbf{Related papers:} Twelfth Thought Experiment {[}S1{]}, double
slit with localized waves {[}S2{]}, earlier two-wave exchange experiment
{[}S3{]}, finite-order resonance {[}S4{]}, exchange weight {[}S5{]},
re-examination of xyztRQ {[}S6{]}, re-run of the earlier experiment
{[}S7{]}

\textbf{Keywords:} coherence, decoherence, observation, wave-packet
collapse, harmonics, localization, scattering, reversible exchange,
periodic closure, internal clock, Lorentz contraction, mass,
quasi-stationary state, particle lifetime

\begin{center}\rule{0.5\linewidth}{0.5pt}\end{center}

\section{Abstract}\label{abstract}

This thought experiment began with a naive question about classifying
wave-packet collapse and decoherence as separate phenomena. Taking
experience with laser holography as a clue, it asks whether coherence is
not an absolute attribute of a single wave but differs according to the
partner wave in the interaction, the harmonic structures of both, their
phases and the readout conditions. A localized wave built from aligned
high harmonics forms a sharp peak, but the spatiotemporal alignment it
needs with another localized wave is extremely strict. Even when
interference with a monochromatic wave can be read, overlap with an
equally localized wave is limited to very narrow conditions.

Next, consider a localized wave colliding with a monochromatic
background wave. Even if the background is initially monochromatic, the
high harmonics of the localized wave can be distributed to the
background by scattering. A new question then arises: even if a particle
acquires some mass-like response through interaction with the
background, does it not at the same time keep losing harmonics and
decay? Looking again at the author's earlier closed two-wave scattering
experiment, localization was not lost in one direction but was
redistributed back and forth between the two waves. A change that looks
locally like dissipation may therefore be a reversible circulation of
the whole system.

If this circulation closes stably, finite periodicity \(U^n=I\) may be
observed \textbf{as a result}. We do not impose \(U^n=I\) first as a law
of interaction in order to explain stability. Rather, we examine which
localized states can survive under the actual interaction, and ask
whether periodic closure, phase locking, quasi-stationarity and
lifetimes appear as results. We further consider whether the temporal
readout of this circulation can correspond to an internal clock delay,
and the spatial readout to the localization width and Lorentz
contraction.

The ``background'' in this paper is not a container with places, times
and a number of waves. It refers to the anonymous part of the total
state that has not yet been distinguished by interaction readouts (§4).
If background and particle are not separate entities, particle species
may be names for conserved quantities of the interaction. This question
is left at the end (§14).

This paper is a \textbf{thought experiment that formulates verification
tasks} arising from a reinterpretation of existing numerical
experiments. It is not a paper that completes derivations of mass, the
Lorentz metric, quantum energy levels or the lifetimes of real
particles.

\begin{center}\rule{0.5\linewidth}{0.5pt}\end{center}

\section{0. Starting Point and Causal
Order}\label{starting-point-and-causal-order}

In the previous Twelfth Thought Experiment {[}S1{]}, without giving
physical names such as time, mass or charge in advance, we asked whether
the apparent physics changes according to the axes from which the same
state and interaction are read. Here we bring that question back to a
concrete two-wave interaction and the localization of harmonics.

The most important caution in the discussion is that \textbf{\(U^n=I\)
is a result, not a cause}.

The order of examination is

\[
\begin{gathered}
\text{redistribution of harmonics and phases by a given interaction}\\
\downarrow\\
\text{stability selection: which localized states can persist}\\
\downarrow\\
\text{results: quasi-stationarity, circulation, finite lifetime}\\
\downarrow\\
\text{read whether periodic closure and integer-ratio phase locking occur}\\
\downarrow\\
\text{read the temporal, spatial and mass-like responses of the same state}
\end{gathered}
\]

The arrows describe the order of exploration; they do not mean that all
the physical correspondences in the lower rows have been confirmed.

\section{1. First Question: For Whom Is Coherence a
Property?}\label{first-question-for-whom-is-coherence-a-property}

The starting point was a question about Hotta's explanation of the
many-worlds interpretation, wave-packet collapse and decoherence. We set
aside the merits of the quantum measurement problem itself. More
basically: in reading what, against what partner, are coherence and
decoherence words?

In laser holography, a tiny vibration rapidly disturbs the interference
image. But if the phase merely shifts by a fixed amount, the fringes
simply move. The case where the phase fluctuates and is averaged during
exposure, and the case where the alignment of high harmonics breaks
down, must be considered separately.

In the author's earlier thought experiment {[}S2{]}, it was shown that
when a localized wave made of odd harmonics is sent through a double
slit, the localized peak interferes as a wave and is read again as a
localized peak, while the alignment condition becomes markedly sensitive
as the number of harmonics increases. From this, the following question
arises.

\begin{quote}
Even for the same wave, interference can be read if the partner is
monochromatic, while if the partner is also a sharply localized wave,
the spatiotemporal window for interaction is extremely narrow. If so,
may we speak of coherence as an absolute property of the target wave
alone?
\end{quote}

Standard wave optics already distinguishes self-coherence and mutual
coherence. This paper does not claim that this distinction was unknown.
The issue is to \textbf{connect it to an interaction that actually
exchanges localization}.

\section{\texorpdfstring{2. A Single Tone and \(10^{60}\) Harmonics:
Sharp Windows in Both Space and
Time}{2. A Single Tone and 10\^{}\{60\} Harmonics: Sharp Windows in Both Space and Time}}\label{a-single-tone-and-1060-harmonics-sharp-windows-in-both-space-and-time}

Consider, for illustration, a one-way travelling wave in which \(N\)
integer harmonics of a fundamental, including even as well as odd
harmonics, are superposed with equal amplitudes and aligned phases.

\[
\Psi_N(x,t)=\frac1N\sum_{h=1}^{N}e^{ih\phi(x,t)},
\qquad \phi(x,t)=k_0x-\omega_0t.
\]

From the finite geometric series, its amplitude is

\[
\Psi_N=e^{i(N+1)\phi/2}
\frac{\sin(N\phi/2)}{N\sin(\phi/2)}.
\]

Within the fundamental period, a very sharp main peak appears near
\(\phi=0\pmod{2\pi}\). The phase width to the first zero is \(2\pi/N\).

\[
\Delta x\sim\frac{\lambda_0}{N},
\qquad
\Delta t\sim\frac{T_0}{N}.
\]

Thus the sharp selectivity due to phase-aligned high harmonics appears
\textbf{not only in the spatial direction but also in the temporal
direction}. However, the one-way travelling wave above is a moving peak
along \(x-ct\); it does not mean that the three spatial directions and
the time direction are localized independently. Implementing
three-dimensional localization requires a separate check of how the
propagation directions are constructed. Temporal localization by phase
alignment of many frequency components is established as laser mode
locking {[}11{]}.

In the discussion we placed the extreme example \(N=10^{60}\). Note that
this is not a measured number of harmonics of a real particle. For
example, if \(\lambda_0\sim10^{26}\,\mathrm m\), then
\(\lambda_0/N\sim10^{-34}\,\mathrm m\), finer even than the atomic
scale. The point is not the numerical value of an absolute length but
that the localization width shrinks roughly in proportion to \(1/N\).

\section{3. Second Question: Why Do Two Localized Waves Interfere
Poorly?}\label{second-question-why-do-two-localized-waves-interfere-poorly}

``The existence of interference'', ``the visibility of interference''
and ``the probability of collision or interaction of localized waves''
are not synonymous. First, consider the normalized mode overlap when
harmonic waves of identical shape overlap with a phase difference
\(\delta\).

\[
C_N(\delta)=\frac1N\sum_{h=1}^N e^{ih\delta},
\qquad
|C_N(\delta)|^2=
\left[\frac{\sin(N\delta/2)}{N\sin(\delta/2)}\right]^2.
\]

For \(N=1\) the magnitude of this overlap does not depend on the phase
difference. For \(N\gg1\), the overlap drops sharply when the two
localized waves are shifted only slightly, and the phase window narrows
to \(O(1/N)\). On the other hand, if the partner is monochromatic,
interference with the common fundamental component can remain. In the
full-mode inner product normalized to equal intensity, its magnitude
becomes \(1/\sqrt N\), so we distinguish the existence of interference
from obtaining a strong signal.

Furthermore, \(\delta=k_0\Delta x-\omega_0\Delta t\). Whether localized
waves actually interact depends not only on the relative phase but also
on the locality of the interaction, the place and time of overlap, and
the coupling law. The question in this model becomes:

\begin{quote}
Before attaching the label ``coherent'' to a wave, should we not state
with which wave, in which temporal and spatial window, and what is read
out?
\end{quote}

\section{4. Third Question: Concretely, What Kind of Wave Is the
``Environment''?}\label{third-question-concretely-what-kind-of-wave-is-the-environment}

\subsection{4.0 How the Word ``Background'' Is
Used}\label{how-the-word-background-is-used}

We first fix how the word ``background'' is used in this and later
sections.

At the starting point of this paper there are no spatial coordinates and
no time coordinate. There are only the total complex state \(Z\) and the
interaction

\[
Z'=F(Z).
\]

Hence ``where the background wave is'' and ``when it interacted'' cannot
be asked before deciding from which relational quantities position and
time are read. The background spacetime is not removed afterwards; it is
not placed in the first place. The index counting computational updates
is not, as it stands, physical time.

Nor can the number of background waves be counted beforehand. The \(N\)
in \(Z=(z_1,\ldots,z_N)\) is the number of components used in the
representation, and the count changes when the basis is changed. Two
waves can be distinguished only when their interaction responses differ.
Likewise, ``with which wave it interacted'' cannot be asked without
names. What we require of the law is covariance under a relabelling
\(P\) of the components,

\[
F(PZ)=P\,F(Z),
\]

and the rule for choosing interaction partners must not depend on hidden
coordinates, orderings or particle labels. What can be observed is
determined afterwards by which relations are conserved and which change
across the interaction. Observation is not a naming performed
afterwards, separately from the interaction.

Therefore the ``background'' in this paper is not a container with
places, times and a count, but a convenient name for the anonymous part
of the total state not yet distinguished by interaction readouts.
Treating the background below as ``a phase-aligned monochromatic complex
wave'' is the simplest illustrative model of this anonymous part, and
does not fix what the background is. The definition of the background
itself will be organized in the next paper.

\subsection{4.1 What, Concretely, Is the
Environment?}\label{what-concretely-is-the-environment}

Standard quantum decoherence theory divides the whole into a subsystem
\(S\) and an environment \(E\), and treats the decay of interference
terms when the unobserved degrees of freedom are removed by a partial
trace {[}1, 2{]}. This mathematical procedure is clear. When
implementing a thought experiment, however, we must specify \textbf{what
the environment is, which degrees of freedom interact, and what state
change remains}.

First, let the background be a phase-aligned monochromatic complex wave.
If the interaction is only a reversible phase rotation of each wave,

\[
A'=Ae^{i\alpha},\qquad B'=Be^{i\beta}.
\]

The amplitudes do not change, while the relational information of
relative phase does. But this alone does not newly generate harmonics,
radiate them outward, or irreversibly lose phase correlations. \textbf{A
change in relational information is different from the occurrence of
decoherence or dissipation.} The spin echo {[}12{]}, in which a signal
that vanished because phases went out of alignment returns by
refocusing, is a classic example of this distinction.

Yet it is also not correct to assert that ``nothing happens if the
background is monochromatic''. The next question changed that view.

\section{5. Fourth Question: What If the Background Is Monochromatic but
the Particle Has
Harmonics?}\label{fourth-question-what-if-the-background-is-monochromatic-but-the-particle-has-harmonics}

The particle side \(P\) has many harmonics, and the background side
\(B\) initially has only the fundamental frequency. If a scattering that
exchanges high-harmonic components between the two waves acts, the
background receives high harmonics.

For illustration, take a linear two-channel exchange for each harmonic
\(h\),

\[
\begin{pmatrix}p_h'\\b_h'\end{pmatrix}
=
\begin{pmatrix}
\cos\theta&-\sin\theta\\
\sin\theta&\cos\theta
\end{pmatrix}
\begin{pmatrix}p_h\\b_h\end{pmatrix}.
\]

If \(b_h(0)=0\) for the high harmonics,

\[
p_h' =p_h\cos\theta,\qquad
b_h'=p_h\sin\theta.
\]

\textbf{That the background was monochromatic does not prevent the
redistribution of harmonics from particle to background.}
Mathematically, no nonlinear operation is needed for the redistribution
itself. Readouts of the localized shape or of the density \(|\Psi|^2\)
are nonlinear quantities, so the apparent change of the waveform can be
complicated. This rotation matrix is illustrative. The implementation of
the earlier experiment is the complex exchange operator \(U_R\) (§6.1),
whose eigenvalues are \(1\) and \(-e^{i\delta}\).

The new question that arose here goes well beyond the initial debate on
decoherence.

\begin{quote}
Even if interaction with the background changes the phase evolution of a
particle and gives it mass-like properties, will the particle not keep
leaking high harmonics to the background and eventually lose its
localization?
\end{quote}

Localized structures that remain without dissipating are already known
in nonlinear field theory. Breathers of the sine-Gordon equation {[}5{]}
and Q-balls {[}6{]}, localized solutions of a complex field whose
internal phase rotates, are examples. These, however, are obtained by
assuming specific potentials or symmetries. In the author's earlier
experiment {[}S3{]} as well, in a closed two-wave exchange without a
radiation term, localization did not dissipate but kept going back and
forth between the two waves. The question of this paper is whether this
non-dissipation, obtained without assuming specific potentials or
symmetries, can be reread as a problem of particle stability and
lifetime.

\section{6. Citing an Earlier Numerical Experiment Again: What Looked
Like Leakage Was a Back-and-Forth
Exchange}\label{citing-an-earlier-numerical-experiment-again-what-looked-like-leakage-was-a-back-and-forth-exchange}

In response to this question, the author presented the closed two-wave
scattering experiment of ``Does a wave packet shrink by observation, or
gather by interaction?'' {[}S3{]} from July 2026. A is initially a broad
low-order wave and B a localized high-order wave; complete transmission,
complete reflection and intermediate scattering were compared. The
earlier experiment reported that, for intermediate scattering,
localization moved between A and B, and that the exchange of
localization continued even when the observer was stopped.

\begin{figure}
\centering
\pandocbounded{\includegraphics[keepaspectratio,alt={Figure 1: Exchange of localization between two waves at R=0.70 (original figure of the earlier numerical experiment)}]{figs/fig01_exchange_R070_waveform_evolution.png}}
\caption{Figure 1: Exchange of localization between two waves at R=0.70
(original figure of the earlier numerical experiment)}
\end{figure}

\textbf{Figure 1.} Original figure of the earlier numerical experiment
{[}S3{]}, which the author presented again during the dialogue of this
thought experiment. The figure shown here was re-output by a faithful
copy of the original program and is pixel-identical to the original
output {[}S7{]}. The collision counts shown are \(0,1,2,3,5,10,20,42\);
the vertical axis is the normalized density readout of each waveform
\texttt{rho\_\allowbreak{}ch\ /\allowbreak{}\ max}, the horizontal axis is \texttt{chi\ /\allowbreak{}\ pi},
and the setting shown is \texttt{R=0.70}. Blue is A and orange is B.
\(R\) is the reflection probability \(|r|^2\) of one collision (in the
original program \(t=e^{i\delta/2}\cos(\delta/2)\),
\(r=-ie^{i\delta/2}\sin(\delta/2)\), \(\delta=2\arcsin\sqrt R\)). The
numbers below are transcribed from the figure legends. The re-run with a
faithful copy of the original program {[}S7{]} confirmed that the same
values are obtained (§6.1).

{\def\LTcaptype{none} % do not increment counter
\begin{longtable}[]{@{}
  >{\raggedleft\arraybackslash}p{(\linewidth - 6\tabcolsep) * \real{0.2500}}
  >{\raggedleft\arraybackslash}p{(\linewidth - 6\tabcolsep) * \real{0.2500}}
  >{\raggedleft\arraybackslash}p{(\linewidth - 6\tabcolsep) * \real{0.2500}}
  >{\raggedleft\arraybackslash}p{(\linewidth - 6\tabcolsep) * \real{0.2500}}@{}}
\toprule\noalign{}
\begin{minipage}[b]{\linewidth}\raggedleft
Collisions
\end{minipage} & \begin{minipage}[b]{\linewidth}\raggedleft
effective order of A \(N_A\)
\end{minipage} & \begin{minipage}[b]{\linewidth}\raggedleft
effective order of B \(N_B\)
\end{minipage} & \begin{minipage}[b]{\linewidth}\raggedleft
sum of shown values
\end{minipage} \\
\midrule\noalign{}
\endhead
\bottomrule\noalign{}
\endlastfoot
0 & 1.0 & 32.0 & 33.0 \\
1 & 10.3 & 22.7 & 33.0 \\
2 & 27.0 & 5.96 & 32.96 \\
3 & 31.1 & 1.87 & 32.97 \\
5 & 2.8 & 30.2 & 33.0 \\
10 & 7.78 & 25.2 & 32.98 \\
20 & 22.2 & 10.8 & 33.0 \\
42 & 16.6 & 16.4 & 33.0 \\
\end{longtable}
}

The effective orders \(N_A,N_B\) are the harmonic orders averaged with
the weights of each wave's harmonic distribution,
\(N_{\rm eff}=\sum_n n\,w_n\).

The author further interpreted the phase progression beyond the figure
as follows: ``at some stage the two waveforms become the same, pass a
position corresponding to 180 degrees of the phase cycle, and at 360
degrees separate back into the original two waves''.

This back-and-forth exchange has already been treated in the author's
later papers {[}S5{]} and {[}S4{]}, as a sweep of the reflectance \(R\)
and as an eigenvalue analysis of the two-channel exchange operator.

What matters here is that \textbf{a temporary decrease of harmonics seen
from the particle side can no longer be interpreted directly as
irreversible dissipation}. That apparent coherence collapses and then
revives is known even in a simple quantum model {[}3{]}.

\subsection{6.1 What the Re-run of the Original Program
Confirmed}\label{what-the-re-run-of-the-original-program-confirmed}

The original program of Figure 1 (the 2026-07-13 preliminary experiment
of {[}S3{]}) was copied without changing a single character and re-run,
and the state at each collision was recorded {[}S7{]}.

\textbf{Control.} The output of the copy was identical to the original
output in all 6163 CSV rows and in the verdict JSON, and the seven
output figures, including Figure 1, were pixel-identical.

\textbf{The normalization does nothing.} The update of the original
program is

\[
a_{k+1}=\mathrm{normalize}(r\,a_k+t\,b_k),\qquad b_{k+1}=\mathrm{normalize}(t\,a_k+r\,b_k).
\]

A (\(q=+1\), identification oscillation \(m=1\)) and B (\(q=-1\),
\(m=2\)) are orthogonal (overlap \(|\langle a,b\rangle|\) at most
\(5.4\times10^{-15}\)), and the deviation from 1 of each squared norm
before normalization was at most \(2.7\times10^{-15}\). The update is
therefore the pure exchange operator
\(U_R=\begin{pmatrix}r&t\\t&r\end{pmatrix}\) itself; removing the
normalization changes \(N_{\rm eff}\) by at most \(3.6\times10^{-14}\).
This recursive map has no dissipation term, and the total power of each
wave stays at 1.

\textbf{The peak height changes.} Figure 1 draws each line scaled to a
maximum of 1. The peak of the density \(\sum_\eta|\psi|^2\), not divided
by its total, is

{\def\LTcaptype{none} % do not increment counter
\begin{longtable}[]{@{}rrr@{}}
\toprule\noalign{}
Collisions & peak of A & peak of B \\
\midrule\noalign{}
\endhead
\bottomrule\noalign{}
\endlastfoot
0 & 0.0039 & 0.1250 \\
1 & 0.0402 & 0.0887 \\
2 & 0.1056 & 0.0233 \\
3 & 0.1216 & 0.0073 \\
5 & 0.0109 & 0.1180 \\
10 & 0.0304 & 0.0985 \\
20 & 0.0867 & 0.0422 \\
42 & 0.0648 & 0.0642 \\
\end{longtable}
}

and varies by a factor of about 30, between about 0.004 and 0.125, while
the total power of each wave stays at 1. When localization moves to the
partner, the wave spreads with its power conserved and its peak goes
down.

\textbf{The sum of effective orders is exactly 33.} For all collisions 0
to 128,

\[
N_A(k)=1+31\sin^2\frac{k\omega}{2},\qquad N_A(k)+N_B(k)=33,\qquad \omega=\pi+2\arcsin\sqrt R,
\]

held (maximum difference from the closed form \(1.5\times10^{-12}\),
from the sum 33 \(7.1\times10^{-14}\)). \(\omega\) is the phase
difference between the two eigenvalues \(1\) and \(-e^{i\delta}\) of
\(U_R\), and \(\sin^2(k\omega/2)\) is the fraction that has moved to the
partner by the \(k\)-th collision. A (order 1) and B (odd harmonics 1 to
63, mean order 32) mix in this fraction, so the sum is conserved. The
fraction moved in the first collision is \(\sin^2(\omega/2)=1-R=0.30\),
indistinguishable from the transmittance \(T\).

\textbf{Reading 180 and 360 degrees.} Read in terms of the exchange
phase \(k\omega\), the two waves are equally divided at
\(k\omega\equiv90^\circ,270^\circ\) (\(N_A=N_B=16.5\)), completely
interchanged at \(k\omega\equiv180^\circ\), and back at
\(k\omega\equiv360^\circ\). The values \(N_A=16.6\), \(N_B=16.4\) at
collision 42 are near the equal division. For \(R=0.70\), however,
\(\omega/2\pi=0.8155\ldots\) is irrational (\(\cos\delta=1-2R=-0.4\) is
not a value that gives \(U^n=I\)), and an exact return does not occur.
Within 128 collisions the closest approach to a return is collision 103,
with exchange phase \(358.55^\circ\) and a fraction \(1.6\times10^{-4}\)
remaining on the partner side. An exact return occurs only when \(R\)
lies on a finite-order root \(R_{n,m}=\cos^2(\pi m/n)\) {[}S4{]}.

\begin{figure}
\centering
\pandocbounded{\includegraphics[keepaspectratio,alt={Figure 2: Pre-normalization norm, effective order and absolute peak height recorded in the R=0.70 re-run}]{figs/fig02_normalization_trace_R070.png}}
\caption{Figure 2: Pre-normalization norm, effective order and absolute
peak height recorded in the R=0.70 re-run}
\end{figure}

\textbf{Figure 2.} Result of recording collisions 0 to 128 with a
faithful copy of the original program of Figure 1 {[}S7{]}. Top: the
amount removed by the normalization at each collision (squared norm
before normalization − 1, at the level of rounding error). Middle:
effective orders \(N_A,N_B\) (solid: original update; dashed: pure
\(U_R\) without normalization; they overlap). Bottom: absolute peak
height of the density not divided by its total. Data:
\texttt{検証\_\allowbreak{}R070交換の正規化と振幅\_\allowbreak{}20261009/\allowbreak{}results/\allowbreak{}trace\_\allowbreak{}renormalized\_\allowbreak{}R070.csv},
\texttt{trace\_\allowbreak{}pure\_\allowbreak{}unitary\_\allowbreak{}R070.csv}, \texttt{summary.json};
recording program:
\texttt{検証\_\allowbreak{}R070交換の正規化と振幅\_\allowbreak{}20261009/\allowbreak{}measure\_\allowbreak{}normalization\_\allowbreak{}trace.py}.

\section{\texorpdfstring{7. Fifth Question: Is \(U^n=I\) Not a Result
Rather Than a
Cause?}{7. Fifth Question: Is U\^{}n=I Not a Result Rather Than a Cause?}}\label{fifth-question-is-uni-not-a-result-rather-than-a-cause}

In the illustrative two-channel rotation model,

\[
S(\theta)^k=
\begin{pmatrix}
\cos(k\theta)&-\sin(k\theta)\\
\sin(k\theta)&\cos(k\theta)
\end{pmatrix}.
\]

If \(n\theta=2\pi m\) (integer \(m\)) for some \(n\),

\[
S^n=I.
\]

So harmonics can move once to the partner and, after several exchanges,
return.

But \textbf{if \(S^n=I\) is imposed as the first assumption, periodic
return is a matter of course}. That explains neither stability nor
quantization. This paper asks the reverse: why does a system driven by a
given exchange law close, or not close?

In particular, we distinguish the following three.

\begin{enumerate}
\def\labelenumi{\arabic{enumi}.}
\tightlist
\item
  \textbf{Apparent recovery:} the normalized density waveform returns to
  a similar shape.
\item
  \textbf{State recovery:} the complex amplitudes of both A and B, all
  harmonics, phases and the pre-normalization scale return to their
  initial values.
\item
  \textbf{Finite order of the operator:} \(U^n=I\) holds for all allowed
  states.
\end{enumerate}

Neither 2 from 1 nor 3 from 2 follows automatically. For example, if
only one initial state lies on a periodic orbit, the whole operator need
not have finite order. Periodic closure is \textbf{to be confirmed as a
computational result}.

In the implementation of the earlier experiment the update is \(U_R\)
itself on the two-dimensional space spanned by the two waves, so state
recovery and finite order of the operator on that space reduce to the
same condition \((-e^{i\delta})^n=1\). \(R=0.70\) does not satisfy it
(§6.1).

\section{8. Sixth Question: As What Is a Non-dissipating Internal
Circulation
Observed?}\label{sixth-question-as-what-is-a-non-dissipating-internal-circulation-observed}

Suppose the particle exchanges harmonics with the background and the
same localized peak can be read repeatedly from outside. Then, even with
no dissipation overall, a \textbf{circulation required for the exchange
with the background} exists inside the particle. The question the author
raised here is:

\begin{quote}
Is what looked like the acquisition of mass an apparent delay of an
internal clock caused by this circulation?
\end{quote}

The ``clock'' in this question is not an ideal clock given a time axis
from outside in advance. It is a period read from the relative phases of
interacting complex waves and the recurrence of the localized peak. For
example, compare the readout phases of a free wave \(F\) and a localized
wave \(P\) with circulation, and examine how

\[
\Delta\varphi(k)=\varphi_P(k)-\varphi_F(k)
\]

accumulates with interaction steps. The recurrence period of the
localized peak, the phase rotation of each mode and the recovery period
of the full state are not necessarily the same.

We must not conclude that ``the clock is necessarily delayed because
there is circulation''. The sign, size and comparison reference of the
delay should be determined from the actual interaction and readout map.

The idea of attaching an internal period to matter goes back to de
Broglie {[}7{]}, and the idea of associating mass with a localized
structure of a field appears in Wheeler's geon {[}10{]}. De Broglie,
however, gave the internal period from a known mass, whereas this paper
asks the reverse direction: reading a mass-like response from internal
circulation.

\section{9. Seventh Question: Does the Spatial Metric Not Change at the
Same
Time?}\label{seventh-question-does-the-spatial-metric-not-change-at-the-same-time}

Next the author added an important correction. Clock delay alone is a
one-sided correspondence with relativity. In special relativity, the
delay of proper clocks and the contraction of length in the direction of
motion are contained in the same Lorentz transformation.

\[
\Delta\tau=\frac{\Delta t}{\gamma},\qquad
L=\frac{L_0}{\gamma},\qquad
\gamma=(1-v^2/c^2)^{-1/2}.
\]

However, entering these formulas into a computer program does not
\textbf{derive} relativity. Building on the observation maps examined in
the Twelfth Thought Experiment {[}S1{]}, this paper proposes to read
separately the temporal phase period and the spatial localization width
from the same harmonic-exchange state, and to check whether they
\textbf{move together by the same function} when the relative motion is
changed.

We do not immediately identify ``the internal clock due to circulation
with the background'' with ``velocity-dependent time dilation''.
Moreover, claiming \textbf{the acquisition of mass} requires, beyond the
change of the clock, correspondence with the momentum--energy relation,
inertial response and so on. At the present stage, mass is a
hypothetical reading name.

\section{10. Eighth Question: Why Phase Locking at Integer
Ratios?}\label{eighth-question-why-phase-locking-at-integer-ratios}

If waves that can exist stably close with a finite period as a result,
integer ratios may appear between several phase rotations. If two
rotations return exactly with a common period \(T\),

\[
\omega_1T=2\pi p,\qquad
\omega_2T=2\pi q,
\quad p,q\in\mathbb Z,
\]

and therefore

\[
\frac{\omega_1}{\omega_2}=\frac pq.
\]

But \textbf{integer ratios must not be given from the start}. Moreover,
in a closed reversible system, quasi-periodic motion with an irrational
ratio can also persist without dissipation. For integer ratios to be
selected as stable states, an actual open path, coupling or stabilizing
mechanism that removes non-aligned states is needed. A precedent for
treating phase locking not as a condition given at the start but as a
result of the dynamics is the synchronization of coupled oscillators
{[}13{]}.

The inference this paper wants to test is:

\begin{quote}
Among waves exchanging harmonics through interaction, there are states
that can circulate and states that pass localized components to paths
that do not return. If the latter are lost and only the former are
observed for a long time, are not integer-ratio phase locking and
discrete quasi-stationary states apparently selected?
\end{quote}

Here ``phase locking'' is not a cause but \textbf{a result that may
appear after stability selection}.

\section{11. Ninth Question: Why Do Particles with Lifetimes
Exist?}\label{ninth-question-why-do-particles-with-lifetimes-exist}

Here the author pointed out that not only stable particles but also
particles with finite lifetimes actually exist. Therefore not only fully
closing states but also incomplete circulations and quasi-stationary
states must be treated in the same framework.

Schematically, we can distinguish

\begin{itemize}
\tightlist
\item
  \textbf{Stable localized states:} even as the exchange between object
  and background continues, the property of being read as a localized
  wave persists.
\item
  \textbf{Metastable localized states:} maintained for a long time, but
  leaking to other degrees of freedom through a slight coupling.
\item
  \textbf{Short-lived states:} redistributing the localized structure in
  a short time to other waves or several particle states.
\end{itemize}

Evaluating a decay that looks irreversible requires a model that
includes whether components that went to external channels return. A
purely closed two-wave exchange alone generally cannot derive
exponential particle lifetimes or decay rates. Bound states in the
continuum, where radiation is suppressed by interference even though
radiating paths exist {[}4{]}, their counterpart under periodic driving
{[}9{]}, and long-lived localized oscillations that keep radiating
slightly (oscillons) {[}8{]} are precedents for quantitatively comparing
the stable and the metastable.

If the survival probability \(P(t)\) of a localized state approximately
follows \(e^{-t/\tau}\), \(\tau\) can be read as a lifetime. Standard
quantum theory has the correspondence \(\Gamma_E\simeq\hbar/\tau\)
between resonance linewidth and lifetime, but this formula is not used
as a starting assumption of this model. The goal is first to see whether
a distribution of lifetimes emerges from the interaction, and then to
compare with known relations.

This also connects to the question of quantum energy levels. However,
\textbf{there is as yet no result that integer-ratio phase locking,
discrete energies and decay lifetimes have been reproduced from the same
numerical model}.

\section{12. Verification Method: Do Not Embed the Conclusion in the
Initial
Conditions}\label{verification-method-do-not-embed-the-conclusion-in-the-initial-conditions}

To turn this thought experiment into computation, the verification needs
the following order.

\subsection{12.1 First Reproduce the Old
Experiment}\label{first-reproduce-the-old-experiment}

Using \textbf{the same interaction matrix, state representation and
readout as the implementation of the time} of the earlier two-wave
exchange experiment {[}S3{]}, reproduce the waveforms at \texttt{R=0.70}
and collision counts \(0,1,2,3,5,10,20,42\). Recheck the controls of
complete transmission and complete reflection. Identify the definition
of the effective-order indicators \(N_A,N_B\) in Figure 1.

\textbf{In this paper, the experiment was re-run with a faithful copy of
the original program {[}S7{]}.} Data and figures all matched the
original output; the definition and closed form of \(N_{\rm eff}\), the
action of the normalization and the absolute peak heights are given in
§6.1. The procedure and package of the re-run are in
\texttt{検証\_\allowbreak{}R070交換の正規化と振幅\_\allowbreak{}20261009/\allowbreak{}README.md},
\texttt{run\_\allowbreak{}all.sh} and \texttt{SHA256SUMS}.

Note that \(R\) in the old experiment is an input value, and its
periodicity is fixed in advance by the eigenvalues of the exchange
operator {[}S4{]}. To verify \(U^n=I\) as a result, a separate model
that determines the mixing angle from the interaction is therefore
needed.

\subsection{12.2 Track the Complex State Itself, Not Only the
Waveform}\label{track-the-complex-state-itself-not-only-the-waveform}

Record the complex coefficients \(p_h(k),b_h(k)\) of each harmonic at
each step, and the difference from the initial state

\[
\varepsilon(k)=
\frac{\|\Psi(k)-\Psi(0)\|}{\|\Psi(0)\|}.
\]

Compute separately the difference with only the global phase removed, so
as not to confuse it with the exact difference. Record
\textbf{separately} candidate indicators corresponding to the norm, the
square closure and the energy. Do not judge ``\(U^n=I\)'' from agreement
of normalized density waveforms alone.

\subsection{12.3 Compare Repeated Exchange with the Same Background and
Leakage to New
Background}\label{compare-repeated-exchange-with-the-same-background-and-leakage-to-new-background}

Compare a closed condition in which reversible exchange is repeated
between the same two waves with an open condition in which components
can be transferred to new background degrees of freedom at each step. In
the latter, keep the state of the new background channels explicitly in
the computational model, and do not casually erase components that
disappeared from the object. Observe leakage to the background,
re-inflow and the lifetime of the localized peak at the same time.

\subsection{12.4 Measure Phase Locking
Afterwards}\label{measure-phase-locking-afterwards}

Without restricting exchange angles or periods to rational numbers,
sweep the initial phases and coupling parameters and measure the phase
advance rates of localized solutions that remain for a long time.
Examine whether long-surviving solutions accumulate near integer ratios,
and whether non-aligned quasi-periodic solutions also survive long.
\textbf{Do not give integer ratios in advance as a setting.}

\subsection{12.5 Read Time and Space Independently and Compare
Afterwards}\label{read-time-and-space-independently-and-compare-afterwards}

For each localized solution, obtain the change of relative phase against
an external reference wave (the internal clock) and the readout of
spatial localization width and direction of motion. Without changing the
interaction, compare afterwards whether the changes with relative motion
are consistent with the Lorentz metric. Do not embed \(\gamma\) from the
start.

\subsection{12.6 Falsifiable Conditions}\label{falsifiable-conditions}

If the following occur, the strong expectations of this thought
experiment are restricted or denied.

\begin{enumerate}
\def\labelenumi{\arabic{enumi}.}
\tightlist
\item
  Even under closed repeated interaction, the full complex state does
  not return and leaks monotonically outward.
\item
  Stable localized states are generally quasi-periodic, and no selection
  by integer-ratio phase locking appears.
\item
  Stable harmonic circulation is found, but no internal clock delay
  appears, or it does not move together with the spatial readout.
\item
  The relation between readable mass-like responses and dissipation,
  circulation and localization width is inconsistent with the features
  of existing real particles.
\end{enumerate}

\section{13. What This Thought Experiment Has Confirmed and Not Yet
Confirmed}\label{what-this-thought-experiment-has-confirmed-and-not-yet-confirmed}

{\def\LTcaptype{none} % do not increment counter
\begin{longtable}[]{@{}
  >{\raggedright\arraybackslash}p{(\linewidth - 2\tabcolsep) * \real{0.5000}}
  >{\raggedright\arraybackslash}p{(\linewidth - 2\tabcolsep) * \real{0.5000}}@{}}
\toprule\noalign{}
\begin{minipage}[b]{\linewidth}\raggedright
Item
\end{minipage} & \begin{minipage}[b]{\linewidth}\raggedright
Status at present
\end{minipage} \\
\midrule\noalign{}
\endhead
\bottomrule\noalign{}
\endlastfoot
Phase alignment of high harmonics creates a sharp localized peak and
narrows the alignment tolerance & Mathematical consequence of the finite
harmonic sum. Theme of earlier work {[}S2{]} \\
Mutual interference properties change with the partner wave and readout
conditions & Can be organized mathematically as wave optics \\
Existing harmonics of a localized wave move to a monochromatic
background & Possible in the two-channel exchange model. Nonlinear
action is not required \\
In intermediate scattering, localization between A and B is
redistributed oscillatorily & Based on the earlier numerical experiment
{[}S3{]} and Figure 1 \\
A closed two-wave exchange has no dissipation; each wave's total power
is conserved while its peak height changes & Confirmed by the re-run of
the original program {[}S7{]}. The normalization acts only at the level
of rounding error \\
The sum of effective orders \(N_A+N_B=33\) is conserved & Confirmed by
the re-run {[}S7{]}. \(N_A=1+31\sin^2(k\omega/2)\) holds for collisions
0--128 to \(10^{-12}\) \\
Equal division at exchange phase \(k\omega=90^\circ\), interchange at
\(180^\circ\), return at \(360^\circ\) & Confirmed by the re-run
{[}S7{]}. \(R=0.70\) is not a finite-order root, so no exact return
occurs \\
\(U^n=I\) arises as a result of the interaction & Central hypothesis.
Finite order of the whole operator is unproven \\
Circulation produces a mass-like delay of the internal clock & Untested
new hypothesis \\
Lorentz contraction is obtained from the same circulation & Untested.
Joint derivation of time and space needed \\
Integer-ratio phase locking, levels and particle lifetimes arise from
stability selection & Untested. Numerical experiments including open
channels needed \\
\end{longtable}
}

\section{14. The Question That Remains Next: Are Particle Species
Reading Names of Conserved
Quantities?}\label{the-question-that-remains-next-are-particle-species-reading-names-of-conserved-quantities}

Treating the background in §4.0 as an anonymous part raises another
question. If background and particle are not separate entities, where do
particle species come from?

Returning to the author's earlier examination of axis signs {[}S6{]},
the signs can be considered not binary but ternary,

\[
\sigma_i\in\{-1,0,+1\}.
\]

\(0\) does not mean that the component is absent but that it is not made
to contribute to the inner product of that readout. Writing the internal
state for illustration as \(v=(t,R,Q_1,Q_2,Q_3)\), a readout candidate
is

\[
\mathcal R_\sigma(v)=\sigma\cdot v=\sum_{i=1}^{5}\sigma_iv_i.
\]

This is an interpretive example treating 3 of the 8 directions as
spatial readouts; it is not a proposal to input named axes \(xyz\) or
\(t,R,Q_i\) into the universal interaction. The formal number of
combinations of three values over five components, \(3^5=243\), does not
mean a number of particles either. It is merely a symbolic count before
considering the anonymity of axes, the dynamics and the equivalence
relations due to observational indistinguishability.

The order considered here is:

\[
\boxed{\begin{array}{c}
\text{universal interaction}
\longrightarrow
\text{conserved / non-conserved relations}
\\
\longrightarrow
\text{distinction by interaction}
\longrightarrow
\text{reading name: particle species}
\end{array}}
\]

Rather than making particle species an input of the interaction, we ask
what is conserved under the same interaction and whether that
conservation relation becomes distinguishable through interaction with
other states.

The first thing to check is whether the existing universal interaction
\(F\) has a sign-conserving property. No new conservation law is added.
Using the existing \(F\) as it is, we examine whether

\[
\boxed{\mathcal I(F(Z))=\mathcal I(Z)}
\]

holds for an invariant \(\mathcal I\) representing the sign of rotation
axes or the relative sense of rotation. If it holds, the interaction
itself has a conservation structure before periodicity or particle
species are assumed.

Further, sign conservation and the stability of a localized structure
may be separate conditions. When a particle decays and changes into
another group of particles, the localized structure breaks, but the
overall conservation relation of rotation signs may be maintained. If
so, a way opens to treat stable and unstable particles in one framework.
The question of lifetimes in §11 is then placed on top of this
distinction.

However, which of charge, spin or particle species this sign corresponds
to must be determined from interaction and observation. Sign
conservation, the correspondence with particle species, and the physical
meaning of the rotation and distribution among 6 directions are, at the
present stage, tasks for verification.

\begin{quote}
Is it not that particle species are not conserved as such, but that
particle species appear distinguishable as a result of reading out the
conserved quantities of the universal interaction?
\end{quote}

\begin{center}\rule{0.5\linewidth}{0.5pt}\end{center}

\section{15. Conclusion: Does Being Able to Exist Stably Come
First?}\label{conclusion-does-being-able-to-exist-stably-come-first}

This dialogue began with the question of how to define decoherence. The
important turning point, however, was reaching the question: \textbf{if
a localized particle's harmonics move to a monochromatic background, why
does the particle not keep dissipating?} Here the meaning of the earlier
two-wave exchange experiment changed.

Passing harmonics to the background in one scattering does not mean that
information is lost from the whole system. In a closed system those
harmonics can return to the particle. If there is periodic circulation,
the particle's localized structure can be maintained; if there is a path
leaking out of the circulation, a quasi-stationary state can have a
lifetime.

Then stability, periodicity, integer-ratio phase locking, internal
clocks, spatial localization, and even mass and particle lifetime,
several reading names, may not be mutually unrelated principles but
different faces of one problem: \textbf{which properties of the same
interaction can exist for a long time}.

But here the order must not be mistaken.

\[
\boxed{\begin{array}{c}
\text{It is not stable because it has a period.}
\\
\text{First there is the interaction; ask what remains stable.}
\\
\text{If periods or integer ratios appear, they are results.}
\end{array}}
\]

If the Twelfth Thought Experiment was an attempt ``not to give names of
physical quantities first'', the Thirteenth Thought Experiment is an
attempt ``\textbf{not to give periodicity conditions or mass first}''.
Common to both is the principle of not pushing properties that could be
read back into the starting point.

\subsection{The Last Question}\label{the-last-question}

\begin{quote}
When a monochromatic background wave and a localized harmonic wave
interact without imposing any special periodicity condition, \textbf{if
we select only the states that can exist stably}, can the period of
those states, the delay of the internal clock, spatial localization,
integer-ratio phase locking and the presence or absence of a finite
lifetime all be read out at the same time?
\end{quote}

\begin{center}\rule{0.5\linewidth}{0.5pt}\end{center}

\section{Reproducing the Figures}\label{reproducing-the-figures}

Figures 1 and 2 are outputs of the program of the earlier experiment.
Running \texttt{検証\_\allowbreak{}R070交換の正規化と振幅\_\allowbreak{}20261009/\allowbreak{}run\_\allowbreak{}all.sh}
performs, in order, the run of the copy identical character by character
to the 0713 original (control, about 5.5 minutes) and the run of the
recording program, and regenerates both figures. The results of the
comparison with the original output are in \texttt{README.md} in the
same folder (Python 3.9.6, numpy 2.0.2, matplotlib 3.9.4, no random
numbers).

{\def\LTcaptype{none} % do not increment counter
\begin{longtable}[]{@{}
  >{\raggedright\arraybackslash}p{(\linewidth - 4\tabcolsep) * \real{0.3333}}
  >{\raggedright\arraybackslash}p{(\linewidth - 4\tabcolsep) * \real{0.3333}}
  >{\raggedright\arraybackslash}p{(\linewidth - 4\tabcolsep) * \real{0.3333}}@{}}
\toprule\noalign{}
\begin{minipage}[b]{\linewidth}\raggedright
Figure
\end{minipage} & \begin{minipage}[b]{\linewidth}\raggedright
File
\end{minipage} & \begin{minipage}[b]{\linewidth}\raggedright
Generated by
\end{minipage} \\
\midrule\noalign{}
\endhead
\bottomrule\noalign{}
\endlastfoot
1 & \texttt{figures/\allowbreak{}fig01\_\allowbreak{}exchange\_\allowbreak{}R070\_\allowbreak{}waveform\_\allowbreak{}evolution.png} &
output
\texttt{exchange\_\allowbreak{}scattering\_\allowbreak{}matrix\_\allowbreak{}R070\_\allowbreak{}waveform\_\allowbreak{}evolution\_\allowbreak{}v1.png}
of
\texttt{検証\_\allowbreak{}R070交換の正規化と振幅\_\allowbreak{}20261009/\allowbreak{}original\_\allowbreak{}copy/\allowbreak{}run\_\allowbreak{}exchange\_\allowbreak{}scattering\_\allowbreak{}matrix\_\allowbreak{}fermionic\_\allowbreak{}localization\_\allowbreak{}transfer\_\allowbreak{}preliminary\_\allowbreak{}v1.py} \\
2 & \texttt{figures/\allowbreak{}fig02\_\allowbreak{}normalization\_\allowbreak{}trace\_\allowbreak{}R070.png} & output
\texttt{results/\allowbreak{}normalization\_\allowbreak{}trace\_\allowbreak{}R070.png} of
\texttt{検証\_\allowbreak{}R070交換の正規化と振幅\_\allowbreak{}20261009/\allowbreak{}measure\_\allowbreak{}normalization\_\allowbreak{}trace.py} \\
\end{longtable}
}

\begin{center}\rule{0.5\linewidth}{0.5pt}\end{center}

\textbf{Note:} This paper is a thought experiment and clearly
distinguishes the earlier numerical experiment on the exchange of
localization by interaction from the newly proposed hypotheses on mass,
the Lorentz metric, levels and lifetimes. The earlier numerical
experiment was re-run with a faithful copy of the original program, and
the output was confirmed to match {[}S7{]}.

\begin{itemize}
\tightlist
\item
  §4.0 and §14 integrate into the main text the content recorded as
  appendix v0.1 from the latter half of the same day's dialogue (the
  author's questions on the place, time, count and interaction partners
  of the background, and the examination of sign conservation of the
  universal interaction). They contain no new numerical experiment,
  proof or establishment of a conservation law.
\item
  Figure 1 is the \textbf{figure of the earlier numerical experiment}
  presented by the author during the dialogue. §6.1 and Figure 2 are
  based on the re-run {[}S7{]}, in which that original program was
  copied without changing a character, its output compared with the
  original, and the states of collisions 0 to 128 recorded. No new
  computation of mass or lifetime was performed.
\item
  Among the formulas, the finite harmonic sum, the real two-channel
  rotation and the overlap are illustrative analytical models. The
  exchange operator of the earlier experiment's implementation was
  identified in §6.1.
\item
  External references {[}3{]}--{[}13{]} were integrated into the main
  text from the list of external prior work compiled the same day
  (\texttt{external\_\allowbreak{}references\_\allowbreak{}thought\_\allowbreak{}experiment\_\allowbreak{}13\_\allowbreak{}ja\_\allowbreak{}v1.0.md}).
  Not all sections of each reference were read closely.
\end{itemize}

\begin{center}\rule{0.5\linewidth}{0.5pt}\end{center}

\section{Record of AI Involvement}\label{record-of-ai-involvement}

\begin{itemize}
\tightlist
\item
  \textbf{ChatGPT} (2026-10-09): drafting of the first version from the
  dialogue with the author (reconstruction in the order in which ideas
  arose; structure, distinction of formulas, explicit statement of
  untested items).
\item
  \textbf{Claude} (2026-10-09 to 10): rewriting §6 as a re-citation of
  the author's own paper; adding self-citations {[}S4{]}--{[}S6{]};
  integrating the appendix into the main text (§4.0, §14); citing
  external references {[}3{]}--{[}13{]} in the main text; correcting
  sentences in §5 and §12.1; re-running the earlier experiment with a
  faithful copy of the original program and recording it ({[}S7{]},
  §6.1, Figure 2); organizing the section structure and format; English
  translation. The posing of the problems, the correction of the main
  causal order, the ideas of stability selection and of mass and
  lifetime, the provision of the existing figure, and design and
  judgement are the author's.
\end{itemize}

\begin{center}\rule{0.5\linewidth}{0.5pt}\end{center}

\section{References}\label{references}

\subsection{This research series}\label{this-research-series}

{[}S1{]} Noriaki Kihara, \textbf{The Twelfth Thought Experiment: From
the Fine-Structure Constant to Observation Maps and Gauge Symmetry},
v1.0 (2026-10-08). DOI: \url{https://doi.org/10.5281/zenodo.23226259} .
Manuscript:
\texttt{thought\_\allowbreak{}experiment\_\allowbreak{}12\_\allowbreak{}observation\_\allowbreak{}mapping\_\allowbreak{}gauge\_\allowbreak{}symmetry\_\allowbreak{}ja\_\allowbreak{}v1.0.md}
(in the Twelfth Thought Experiment folder under the same parent folder).

{[}S2{]} Noriaki Kihara, \textbf{The mystery of the double slit: neither
observation nor ``wave-packet collapse'' was needed? --- A particle-like
lump of waves interferes as a wave and appears as a particle} (in
Japanese), note (2026-06-29).
\url{https://note.com/kiharanoriaki/n/n65be6bf06c9b} . DOI of the
localized-wave version of the original study:
\url{https://doi.org/10.5281/zenodo.21035831} .

{[}S3{]} Noriaki Kihara, \textbf{Does a wave packet shrink by
observation, or gather by interaction?} (in Japanese), note
(2026-07-13). \url{https://note.com/kiharanoriaki/n/nbc6649e30af3} . Concept
DOI of the original study: \url{https://doi.org/10.5281/zenodo.21333766} ;
Version DOI: \url{https://doi.org/10.5281/zenodo.21333768} .

{[}S4{]} Noriaki Kihara, \textbf{Discovery of finite-order resonances in
repeated exchange scattering --- identifying the cause of peaks near the
fine-structure constant 137 and 128 and a reproducible wave-packet
mathematical model} (in Japanese), v1 (2026-07-18). Concept DOI:
\url{https://doi.org/10.5281/zenodo.21421366} ; Version DOI:
\url{https://doi.org/10.5281/zenodo.21421367} .

{[}S5{]} Noriaki Kihara, \textbf{Selection of the exchange weight
G\_R=1-R and numerical experiments on candidates corresponding to the
fine-structure constant} (in Japanese), v1 (2026-07-15). Concept DOI:
\url{https://doi.org/10.5281/zenodo.21396760} ; Version DOI:
\url{https://doi.org/10.5281/zenodo.21396761} .

{[}S6{]} Noriaki Kihara, \textbf{Re-examination of the six-dimensional
encoding xyztRQ --- thought-experiment notes toward the next paper} (in
Japanese), v4 (2026-04-30). Concept DOI:
\url{https://doi.org/10.5281/zenodo.19902677} ; Version DOI:
\url{https://doi.org/10.5281/zenodo.19904714} .

{[}S7{]} Noriaki Kihara, \textbf{Verification: normalization and
amplitude in the R=0.70 recursive exchange} (2026-10-09, in the same
folder as this paper). \texttt{検証\_\allowbreak{}R070交換の正規化と振幅\_\allowbreak{}20261009/\allowbreak{}}:
faithful copy of the original program
\texttt{original\_\allowbreak{}copy/\allowbreak{}run\_\allowbreak{}exchange\_\allowbreak{}scattering\_\allowbreak{}matrix\_\allowbreak{}fermionic\_\allowbreak{}localization\_\allowbreak{}transfer\_\allowbreak{}preliminary\_\allowbreak{}v1.py}
and its output (control), recording program
\texttt{measure\_\allowbreak{}normalization\_\allowbreak{}trace.py}, data
\texttt{results/\allowbreak{}trace\_\allowbreak{}renormalized\_\allowbreak{}R070.csv},
\texttt{trace\_\allowbreak{}pure\_\allowbreak{}unitary\_\allowbreak{}R070.csv},
\texttt{control\_\allowbreak{}vs\_\allowbreak{}original\_\allowbreak{}snapshots.csv}, \texttt{summary.json},
figure \texttt{results/\allowbreak{}normalization\_\allowbreak{}trace\_\allowbreak{}R070.png},
\texttt{README.md}, \texttt{run\_\allowbreak{}all.sh}, \texttt{SHA256SUMS}. Included
in the Zenodo record of this paper (10.5281/zenodo.23266251) as
\texttt{verification\_\allowbreak{}R070\_\allowbreak{}exchange\_\allowbreak{}20261009.zip}.

\subsection{External references}\label{external-references}

{[}1{]} W. H. Zurek, \emph{Decoherence, einselection, and the quantum
origins of the classical}, Reviews of Modern Physics \textbf{75},
715--775 (2003). \url{https://doi.org/10.1103/RevModPhys.75.715} .

{[}2{]} M. Schlosshauer, \emph{Decoherence, the measurement problem, and
interpretations of quantum mechanics}, Reviews of Modern Physics
\textbf{76}, 1267--1305 (2005).
\url{https://doi.org/10.1103/RevModPhys.76.1267} .

{[}3{]} J. H. Eberly, N. B. Narozhny, J. J. Sánchez-Mondragón,
\emph{Periodic spontaneous collapse and revival in a simple quantum
model}, Physical Review Letters \textbf{44}, 1323 (1980).
\url{https://doi.org/10.1103/PhysRevLett.44.1323} .

{[}4{]} H. Friedrich, D. Wintgen, \emph{Interfering resonances and bound
states in the continuum}, Physical Review A \textbf{32}, 3231--3242
(1985). \url{https://doi.org/10.1103/PhysRevA.32.3231} .

{[}5{]} K. Maki, H. Takayama, \emph{Quantum-statistical mechanics of
extended objects. II. Breathers in the sine-Gordon system}, Physical
Review B \textbf{20}, 5002 (1979).
\url{https://doi.org/10.1103/PhysRevB.20.5002} .

{[}6{]} S. Coleman, \emph{Q-balls}, Nuclear Physics B \textbf{262},
263--283 (1985). \url{https://doi.org/10.1016/0550-3213(85)90286-X} .

{[}7{]} L. de Broglie, \emph{Waves and quanta}, Nature \textbf{112}, 540
(1923). \url{https://doi.org/10.1038/112540a0} .

{[}8{]} G. Fodor, \emph{A review on radiation of oscillons and
oscillatons}, arXiv:1911.03340 (2019).
\url{https://doi.org/10.48550/arXiv.1911.03340} .

{[}9{]} S. Longhi, G. Della Valle, \emph{Floquet bound states in the
continuum}, Scientific Reports \textbf{3}, 2219 (2013).
\url{https://doi.org/10.1038/srep02219} .

{[}10{]} J. A. Wheeler, \emph{Geons}, Physical Review \textbf{97}, 511
(1955). \url{https://doi.org/10.1103/PhysRev.97.511} .

{[}11{]} H. A. Haus, \emph{Mode-locking of lasers}, IEEE Journal of
Selected Topics in Quantum Electronics \textbf{6}, 1173--1185 (2000).
\url{https://doi.org/10.1109/2944.902165} .

{[}12{]} E. L. Hahn, \emph{Spin echoes}, Physical Review \textbf{80},
580--594 (1950). \url{https://doi.org/10.1103/PhysRev.80.580} .

{[}13{]} Y. Kuramoto, \emph{Self-entrainment of a population of coupled
non-linear oscillators}, Lecture Notes in Physics \textbf{39}, 420--422
(1975). \url{https://doi.org/10.1007/BFb0013365} .

\end{document}
